% ============================================================
%  RAPPORT DE STAGE D'ÉTÉ — SAHALI
%  Institut National des Sciences Appliquées et de Technologie (INSAT)
%  Compile with: xelatex main.tex  (run twice for the table of contents)
%  Requires an Arabic font installed — this file uses "Amiri" (a free,
%  installed on your system, replace \newfontfamily\arabicfont{Amiri}
%  below with any Arabic font you have (e.g. "Noto Naskh Arabic",
%  "Scheherazade New").
% ============================================================
\documentclass[12pt, a4paper]{report}

% ---- Fonts and languages (XeLaTeX) ----
\usepackage{fontspec}
\usepackage{polyglossia}
\setmainlanguage{english}
\setotherlanguage{french}
\setotherlanguage{arabic}
\newfontfamily\arabicfont[Script=Arabic]{Amiri}

% ---- Layout ----
\usepackage[top=2.5cm, bottom=2.5cm, left=3cm, right=3cm]{geometry}
\usepackage{setspace}
\onehalfspacing

% ---- Graphics and colour ----
\usepackage[dvipsnames]{xcolor}
\usepackage{graphicx}
\usepackage{tikz}
\usetikzlibrary{shapes.geometric, arrows.meta, positioning, fit, backgrounds}

% ---- Tables ----
\usepackage{booktabs}
\usepackage{tabularx}
\usepackage{array}
\usepackage{multirow}

% ---- Lists ----
\usepackage{enumitem}

% ---- Code listings ----
\usepackage{listings}
\definecolor{codebg}{HTML}{F5F7FF}
\definecolor{codekw}{HTML}{0038AF}
\definecolor{codestr}{HTML}{15803D}
\definecolor{codecomment}{HTML}{64748B}
\lstdefinestyle{code}{
  backgroundcolor=\color{codebg},
  basicstyle=\ttfamily\footnotesize,
  keywordstyle=\color{codekw}\bfseries,
  stringstyle=\color{codestr},
  commentstyle=\color{codecomment}\itshape,
  breaklines=true,
  frame=single,
  framerule=0.3pt,
  rulecolor=\color{gray!40},
  numbers=left,
  numberstyle=\tiny\color{gray},
  tabsize=2,
  showstringspaces=false,
}
\lstset{style=code}

% ---- Headers/footers ----
\usepackage{fancyhdr}
\pagestyle{fancy}
\fancyhf{}
\fancyhead[L]{\small\textcolor{gray}{Sahali --- Summer Internship Report}}
\fancyhead[R]{\small\textcolor{gray}{\leftmark}}
\fancyfoot[C]{\thepage}
\renewcommand{\headrulewidth}{0.4pt}

% ---- Captions ----
\usepackage[font=small, labelfont=bf, labelsep=colon]{caption}

% ---- Hyperlinks (load last) ----
\usepackage[colorlinks=true, linkcolor=NavyBlue, citecolor=NavyBlue, urlcolor=NavyBlue]{hyperref}

% ---- Custom colours ----
\definecolor{primary}{HTML}{0038AF}
\definecolor{amber}{HTML}{F59E0B}
\definecolor{dark}{HTML}{0D1B35}
\definecolor{muted}{HTML}{64748B}

\begin{document}

% ============================================================
% PAGE DE GARDE
% ============================================================
\begin{titlepage}
  \centering
  \begin{french}
  \begin{minipage}{0.48\textwidth}
    \centering
    RÉPUBLIQUE TUNISIENNE\\
    Ministère de l'Enseignement Supérieur\\
    et de la Recherche Scientifique
  \end{minipage}
  \hfill
  \begin{minipage}{0.48\textwidth}
    \centering
    \textit{[Emplacement du logo INSAT]}\\[2pt]
    Université de Carthage
  \end{minipage}

  \vspace{0.6cm}
  {\Large\bfseries Institut National des Sciences Appliquées et de Technologie}\\[2pt]
  \color{primary}\rule{\textwidth}{1.5pt}
  \end{french}

  \vspace{2cm}

  {\color{muted}\large Rapport de Stage d'Été}

  \vspace{0.8cm}
  {\Huge\bfseries\color{dark} Sahali}

  \vspace{0.3cm}
  {\large\itshape\color{muted} A Civic Reporting Platform for Tunisian Municipalities}

  \vspace{0.5cm}
  {\color{amber}\rule{0.6\textwidth}{1pt}}

  \vspace{2cm}

  \begin{minipage}{0.45\textwidth}
    \begin{flushleft}
      \large
      \textbf{Prepared by:}\\
      [NOM DE L'ÉTUDIANT]
    \end{flushleft}
  \end{minipage}
  \hfill
  \begin{minipage}{0.45\textwidth}
    \begin{flushright}
      \large
      \textbf{Supervised by:}\\
      [NOM DE L'ENCADRANTE]
    \end{flushright}
  \end{minipage}

  \vspace{1cm}
  {\large Host organization: \textbf{Nadas Group}}

  \vfill

  {\large\bfseries Academic Year 2025 -- 2026}

  \vspace{0.5cm}
  {\color{primary}\rule{\textwidth}{1.5pt}}
\end{titlepage}

% ============================================================
% REMERCIEMENTS
% ============================================================
\chapter*{Acknowledgements}
\addcontentsline{toc}{chapter}{Acknowledgements}

I would like to express my sincere gratitude to everyone who supported me throughout this summer internship.

First, my thanks go to my supervisor, [NOM DE L'ENCADRANTE], for her guidance, availability, and constructive feedback at every stage of this project, from the initial architectural decisions to the final security hardening pass.

I would also like to thank Nadas Group for hosting this internship and providing the environment and trust needed to take the Sahali project from an early prototype to a working, tested, and continuously-integrated platform.

My gratitude extends to the faculty of the Institut National des Sciences Appliquées et de Technologie (INSAT), whose teaching in software engineering, databases, and system architecture provided the foundation this project was built on.

Finally, I thank my family and friends for their patience and encouragement throughout this internship.

% ============================================================
% RÉSUMÉ (FR / EN / AR)
% ============================================================
\chapter*{Abstract}
\addcontentsline{toc}{chapter}{Abstract}

\begin{french}
\section*{Résumé (Français)}
Les problèmes d'infrastructure urbaine en Tunisie --- nids-de-poule, éclairage public défaillant, fuites d'eau, dépôts sauvages --- sont aujourd'hui signalés par des canaux fragmentés (appels téléphoniques, visites en personne, réseaux sociaux) qui ne sont ni suivis, ni tracés, ni routés vers le bon service municipal. Ce rapport présente \textbf{Sahali}, une plateforme civique complète développée durant un stage d'été chez Nadas Group, qui structure ce processus en un pipeline unique : signalement, routage automatique vers la municipalité la plus proche, affectation à une équipe terrain, résolution, et retour au citoyen.

La plateforme se compose d'une application mobile Flutter trilingue (arabe, français, anglais) destinée aux citoyens, d'un tableau de bord web React/TypeScript destiné aux équipes municipales, et d'une API backend FastAPI qui orchestre l'authentification, la géolocalisation, les notifications multi-canal et le stockage des données. Une part importante du stage a été consacrée à un audit de sécurité complet (rotation de clés compromises, durcissement de l'authentification, limitation de débit, gestion atomique des jetons) et à la mise en place d'un pipeline d'intégration et de déploiement continus (CI/CD) avec une suite de tests automatisés.

Ce rapport détaille la démarche méthodologique, l'architecture retenue, les choix techniques, les difficultés rencontrées et leurs solutions, ainsi que les résultats des tests de validation.

\textbf{Mots-clés :} plateforme civique, FastAPI, Flutter, PostgreSQL/PostGIS, sécurité applicative, CI/CD, signalement citoyen.
\end{french}

\bigskip

\section*{Abstract (English)}
Urban infrastructure issues in Tunisia --- potholes, faulty streetlights, water leaks, illegal waste dumping --- are today reported through fragmented channels (phone calls, in-person visits, social media posts) that are neither tracked nor routed to the right municipal department. This report presents \textbf{Sahali}, a full civic reporting platform built during a summer internship at Nadas Group, which structures this process into a single pipeline: report, automatic routing to the nearest municipality, assignment to a field team, resolution, and feedback to the citizen.

The platform consists of a trilingual (Arabic, French, English) Flutter mobile application for citizens, a React/TypeScript web dashboard for municipal staff, and a FastAPI backend that orchestrates authentication, geolocation, multi-channel notifications, and data storage. A significant part of the internship was dedicated to a full security audit (rotating a compromised signing key, hardening authentication, rate limiting, atomic token handling) and to building a continuous integration/continuous deployment (CI/CD) pipeline with an automated test suite.

This report details the methodology followed, the chosen architecture, the technical decisions made, the difficulties encountered and how they were resolved, and the results of the validation tests.

\textbf{Keywords:} civic platform, FastAPI, Flutter, PostgreSQL/PostGIS, application security, CI/CD, citizen reporting.

\bigskip

\begin{arabic}
\section*{ملخص (بالعربية)}
اليوم، يتمّ الإبلاغ عن مشاكل البنية التحتية الحضرية في تونس -- مثل الحفر في الطرقات، وأعطال الإنارة العمومية، وتسرّبات المياه، والنفايات العشوائية -- عبر قنوات مبعثرة (مكالمات هاتفية، زيارات شخصية، منشورات على وسائل التواصل الاجتماعي) لا تُتابَع ولا توجَّه إلى المصلحة البلدية المختصة. يقدّم هذا التقرير منصّة \textbf{سهلي}، وهي منصّة مدنية متكاملة تم تطويرها خلال تربّص صيفي لدى مجموعة Nadas Group، تنظّم هذه العملية ضمن مسار واحد: الإبلاغ، ثمّ التوجيه الآلي إلى أقرب بلدية، ثمّ الإسناد إلى فريق ميداني، فالمعالجة، وأخيرًا إعلام المواطن بالنتيجة.

تتكوّن المنصّة من تطبيق هاتفي بلغة فلَتَر (Flutter) ثلاثي اللغة (عربية، فرنسية، إنڤليزية) موجّه للمواطنين، ولوحة تحكّم على الويب بلغتَي React وTypeScript موجّهة لأعوان البلديات، وواجهة برمجية خلفية بإطار FastAPI تتولّى المصادقة، وتحديد الموقع الجغرافي، والإشعارات متعدّدة القنوات، وتخزين البيانات. خُصِّص جزء مهمّ من التربّص لمراجعة أمنية شاملة (تجديد مفتاح توقيع مُخترَق، تعزيز آليات المصادقة، تحديد معدّل الطلبات، معالجة ذرّية للرموز المميّزة) ولبناء سلسلة تكامل ونشر مستمرّين (CI/CD) مع مجموعة اختبارات آلية.

يعرض هذا التقرير المنهجية المتّبعة، والهندسة المعتمدة، والخيارات التقنية، والصعوبات التي ووجهت وكيفية تجاوزها، ونتائج اختبارات التحقّق.

\textbf{الكلمات المفتاحية:} منصّة مدنية، FastAPI، Flutter، PostgreSQL/PostGIS، أمن التطبيقات، CI/CD، الإبلاغ المواطني.
\end{arabic}

% ============================================================
% TABLE DES MATIÈRES
% ============================================================
\newpage
\tableofcontents
\listoffigures
\listoftables
\newpage

% ============================================================
% INTRODUCTION GÉNÉRALE
% ============================================================
\chapter*{General Introduction}
\addcontentsline{toc}{chapter}{General Introduction}

Local governments are, in practice, judged less by policy documents than by whether a broken streetlight gets fixed. In Tunisia, as in many countries, the channel between a citizen who notices a problem and the municipal department able to fix it is informal and lossy: a phone call that nobody logs, a Facebook post that nobody routes, a visit to a counter during limited opening hours. The citizen rarely learns what happened to their report, and the municipality has no aggregate view of what is failing across its territory.

This report documents the design and development of \textbf{Sahali}, a civic reporting platform built to close that loop, carried out as a summer internship at Nadas Group. The work spanned the full stack of a modern web/mobile product: a trilingual citizen-facing mobile application, a municipal staff dashboard, a REST API backend with a geospatial database, and --- once the functional core was working --- a full security audit and the construction of a continuous integration and deployment pipeline. The report plan is detailed at the end of Chapter 1.

% ============================================================
% CHAPITRES
% ============================================================
\chapter{General Context}
\label{chap:cadre-general}

\section{General Context}

Nadas Group hosted this summer internship to build \textbf{Sahali} (\textarabic{سهلي}, ``easy'' / ``made simple'' in Tunisian Arabic), a civic technology platform aimed at Tunisian municipalities. The name reflects the product goal directly: reporting a civic problem should be as easy as taking a photo.

Civic reporting --- the practice of letting residents flag non-emergency infrastructure problems to their local government --- is a well-established category internationally (Section~\ref{sec:etat-art} reviews existing solutions such as FixMyStreet and 311-style systems). What is comparatively underserved is a solution designed around Tunisia's specific constraints: trilingual usage (Arabic, French, and increasingly English among younger users), a fragmented municipal landscape where many communes have limited IT budgets, and address data (via OpenStreetMap) that is frequently incomplete or only partially translated between French and Arabic.

\section{Motivation and Problem Statement}

Before Sahali, a citizen who noticed a pothole, a broken streetlight, or an overflowing waste container had no structured way to report it. In practice, reports happened through phone calls that were not logged, in-person visits during limited office hours, or informal social media posts --- none of which were tracked, routed to the correct department, or followed up on. From the municipality's side, this meant no aggregate, geolocated view of what was failing across the territory, no service-level accountability, and no way to measure resolution performance over time.

This translates into three concrete problems that motivated the project:

\begin{enumerate}[leftmargin=*]
  \item \textbf{For the citizen}: no single, traceable channel to report an issue and follow its resolution.
  \item \textbf{For the municipality}: no centralized tool to triage, assign, and track field interventions.
  \item \textbf{For the citizen--municipality relationship}: no transparency, which erodes trust over time.
\end{enumerate}

\section{Project Objectives}

The internship's objective was to design and build a working, secure, end-to-end platform addressing the three problems above. Concretely, this meant:

\begin{itemize}[leftmargin=*]
  \item A mobile application that lets a citizen submit a categorized, geolocated, photographed report in under two minutes, in their language of choice (Arabic, French, or English).
  \item An automatic routing mechanism that assigns each report to the nearest municipality without manual sorting.
  \item A web dashboard that gives municipal staff (administrators, supervisors, field agents, analysts) a live, filterable view of reports and a way to assign and track field interventions.
  \item A resolution feedback loop: the citizen is notified at each status change and can see the full history of their own report.
  \item A backend built to a production-grade security standard --- not an afterthought bolted on at the end, but audited and hardened as a dedicated phase of the internship.
  \item An automated testing and deployment pipeline, so that regressions are caught before they reach production rather than discovered by users.
\end{itemize}

Deliberately out of scope for this internship (documented as future work in Chapter~\ref{chap:conclusion}): a dedicated AI microservice for automatic report classification and duplicate detection, and formal load testing under realistic concurrent traffic.

\section{Methodology}

The project followed an \textbf{iterative, feedback-driven methodology} rather than a rigid upfront specification. Work was organized in short cycles, each producing a working increment that was immediately exercised end-to-end (via the running application, or via direct API calls) before moving to the next increment --- a lightweight, Kanban-like flow rather than fixed-length Scrum sprints, which suited a one-person internship team better than ceremony-heavy Scrum.

A distinguishing feature of the methodology was its \textbf{audit-driven hardening phase}: once the functional platform was working, a dedicated pass was run specifically to find and fix security issues, structured as:

\begin{enumerate}[leftmargin=*]
  \item \textbf{Audit} --- a systematic review of the codebase across the backend, mobile app, and dashboard, producing a severity-ranked list of findings (critical / high / medium / low).
  \item \textbf{Fix} --- addressing findings in severity order, critical first.
  \item \textbf{Verify} --- for every fix, a direct, reproducible test (a scripted API call sequence, or an automated test) confirming the fix actually changes behavior, not just that the code compiles.
  \item \textbf{Review} --- an independent automated code-review pass (CodeRabbit, integrated with the project's GitHub pull requests) that caught several issues the manual audit had missed, most notably race conditions in concurrent token handling (Section~\ref{sec:difficultes}).
\end{enumerate}

This audit/fix/verify/review loop was then repeated for the continuous integration and deployment (CI/CD) work described in Chapter~\ref{chap:realisation}, and again surfaced real, previously-undetected issues (a broken mobile smoke test, several outdated dependencies with known CVEs), reinforcing the value of treating verification as a first-class part of the methodology rather than a final checkbox.

\section{Project Planning}

Table~\ref{tab:planning} summarizes the internship's timeline, reconstructed retrospectively from the project's development history rather than fixed in advance --- consistent with the iterative methodology described above.

\begin{table}[h!]
\centering
\begin{tabularx}{\textwidth}{llX}
\toprule
\textbf{Phase} & \textbf{Approx.\ duration} & \textbf{Focus} \\
\midrule
1 --- Core platform      & Weeks 1--4   & Backend API foundation, database schema, mobile app skeleton, dashboard skeleton, authentication. \\
2 --- Feature build-out  & Weeks 5--7   & Report submission flow, geocoding and municipality routing, dashboard modules (map, teams, statistics), field-agent mobile mode, multi-media resolution reports. \\
3 --- Security audit \& hardening & Weeks 8--10 & Three-phase severity-ranked remediation: critical (key rotation, auth gating, session revocation), high (rate limiting, headers, CORS), medium/low (dependency hygiene, dead-code cleanup). \\
4 --- CI/CD pipeline     & Weeks 11--12 & Automated test suite, GitHub Actions workflows (lint, dependency audit, Docker build check, mobile release automation), dependency vulnerability remediation. \\
\bottomrule
\end{tabularx}
\caption{Internship timeline by phase}
\label{tab:planning}
\end{table}

\section{Report Plan}

The remainder of this report is organized as follows. \textbf{Chapter~\ref{chap:etude}} covers the study and specification phase: a review of related civic-reporting solutions, the core technical concepts the project relies on, and the functional and non-functional requirements. \textbf{Chapter~\ref{chap:architecture}} details the system architecture and design, first at a technology-agnostic level and then with the concrete technology stack. \textbf{Chapter~\ref{chap:realisation}} describes the implementation: the development environment, key algorithms, and the difficulties encountered along with how they were resolved. \textbf{Chapter~\ref{chap:tests}} presents the testing methodology and results. \textbf{Chapter~\ref{chap:conclusion}} concludes with a project assessment and future perspectives.

\chapter{Study and Specification Phase}
\label{chap:etude}

\section{State of the Art and Study of Existing Solutions}
\label{sec:etat-art}

\subsection{Existing Solutions}

Civic issue-reporting platforms are an established category internationally, and reviewing them clarified which patterns to reuse and which gaps to fill for a Tunisia-specific product.

\begin{itemize}[leftmargin=*]
  \item \textbf{FixMyStreet}~\cite{fixmystreet} (UK, mySociety) --- open-source, one of the earliest and most widely deployed civic reporting platforms. Strong on report-to-map visualization and category-based routing; weaker on native mobile experience and not designed for multilingual/RTL usage.
  \item \textbf{311 systems}~\cite{nyc311} (United States --- e.g.\ NYC 311, SeeClickFix) --- large-scale municipal service-request platforms, typically covering more than just infrastructure (they also handle service requests such as noise complaints). They demonstrate the value of a unified tracking code and public status API, both of which Sahali adopts.
  \item \textbf{E-government municipal portals in the MENA region} --- generally administrative-form-oriented (permits, certificates) rather than reactive infrastructure reporting, and rarely offer a native mobile app with offline support.
\end{itemize}

\subsection{Technologies, Methods, and Tools in the Domain}

Across the solutions above, a common technical pattern emerges: a mobile or web front end for citizens, a relational or geospatial database for report storage, a staff-facing dashboard for triage, and a notification layer (SMS/email/push) for status updates. Increasingly, platforms in this space add a machine-learning layer for automatic category classification and duplicate detection --- a pattern Sahali's design anticipates (see the \texttt{ai\_category\_id} and \texttt{is\_duplicate} fields in the data model, Section~\ref{sec:class-diagram}) even though the model itself is future work.

\subsection{Limitations of Existing Solutions}

None of the reviewed solutions fully addresses the constraints this project needed to satisfy:

\begin{itemize}[leftmargin=*]
  \item Limited or no Arabic/RTL support, which is essential for a Tunisian citizen-facing product.
  \item Not designed around Tunisia's municipal structure (multi-tenant by commune, each with its own staff and service-level targets).
  \item Address data in the region is frequently incomplete on OpenStreetMap and available only in one of French or Arabic per point, which existing platforms do not handle explicitly.
  \item Several are either closed-source with licensing costs unsuited to a pilot-stage municipal budget, or open-source but architected for a single large deployment rather than a multi-tenant one.

\end{itemize}

\section{Core Concepts}

This section defines three concepts referenced throughout the report.

\subsection{CI/CD}

\textbf{Continuous Integration (CI)} is the practice of automatically building and testing every code change --- in this project, on every push, via GitHub Actions --- so that regressions are caught within minutes rather than discovered later in production. \textbf{Continuous Deployment/Delivery (CD)} extends this to automating the release process itself: in this project, a signed Android release build is produced automatically from a version tag rather than built by hand on a developer's machine. Together, CI/CD replace manual, error-prone release steps with a repeatable, auditable pipeline (detailed in Section~\ref{sec:cicd}).

\subsection{Cloud Hosting}

\textbf{Cloud hosting} means running infrastructure (servers, databases, storage) on a third-party provider's platform rather than on physical hardware the team owns. Sahali uses three managed cloud services rather than self-hosted infrastructure: \textbf{Render} for the backend API (Docker-based deployment), \textbf{Vercel} for the dashboard (static site hosting with automatic deployment from Git), and \textbf{Supabase} for the production database and file storage (managed PostgreSQL with the PostGIS extension, plus an S3-compatible storage API). This choice removed the need to provision or maintain servers during the internship, at the cost of some platform-specific configuration (e.g.\ passing cryptographic keys as base64-encoded environment variables rather than files).

\subsection{LLM-Assisted Development}

A \textbf{Large Language Model (LLM)} is a machine learning model trained to understand and generate text (and, in coding-assistant contexts, source code) from natural-language prompts. This project was developed with substantial assistance from an LLM-based coding agent (Claude Code), used interactively throughout: for writing and reviewing code, for running and interpreting test suites, for diagnosing build and deployment failures, and for the systematic security audit described in Chapter~\ref{chap:realisation}. This is disclosed here both for transparency and because it directly shaped the iterative, verify-every-step methodology described in Section~\ref{chap:cadre-general} --- an LLM agent can make changes quickly, which raises rather than lowers the importance of verifying each change actually works before moving to the next one.

\section{Specification}

\subsection{Identification of Actors}

\begin{table}[h!]
\centering
\begin{tabularx}{\textwidth}{lX}
\toprule
\textbf{Actor} & \textbf{Role} \\
\midrule
Citizen        & Submits reports (registered or anonymous), tracks their status, receives notifications. \\
Field Agent    & Municipal staff assigned to specific reports; updates status and files a resolution report on completion. \\
Analyst        & Reviews and reclassifies incoming reports, monitors statistics; no administrative rights. \\
Supervisor     & Assigns reports to one or more field agents; manages teams for their municipality. \\
Admin          & Full platform control: municipalities, staff accounts, category/SLA configuration, cross-municipality visibility. \\
\bottomrule
\end{tabularx}
\caption{System actors}
\label{tab:actors}
\end{table}

\subsection{Functional Requirements}

\begin{itemize}[leftmargin=*]
  \item \textbf{Authenticate users} --- registration and login by phone (OTP) or email/password, for both citizens and staff, with role-based access control.
  \item \textbf{Submit a report} --- category, photo, GPS location (with manual correction), optional description; works for both registered and anonymous citizens.
  \item \textbf{Automatic routing} --- match each report to the nearest municipality by geographic distance, with no manual sorting step.
  \item \textbf{Track a report} --- citizens see their own reports' full status history; a public tracking-code lookup works without an account.
  \item \textbf{Triage and assign} --- municipal staff can reclassify a report's category and assign it to one or more field agents.
  \item \textbf{Update report status} --- along the state machine \texttt{submitted → received → under\_review → in\_progress → resolved} (or \texttt{rejected} at any point), enforced server-side.
  \item \textbf{File a resolution report} --- field agents attach up to five photos, an optional video, an optional voice note, and a mandatory written comment when closing a report.
  \item \textbf{Notify} --- push, SMS, and email notifications dispatched on every status change, plus a live in-dashboard event stream for staff.
  \item \textbf{View statistics} --- aggregate counts, resolution-time averages, and per-municipality performance figures on the dashboard.
  \item \textbf{Manage staff and municipalities} --- administrators can create/update/deactivate staff accounts and municipality records.
\end{itemize}

\subsection{Non-Functional Requirements}

\begin{table}[h!]
\centering
\begin{tabularx}{\textwidth}{lX}
\toprule
\textbf{Requirement} & \textbf{Target} \\
\midrule
Multilingual        & Mobile app: Arabic (RTL), French, English. Dashboard: French, Arabic. \\
Responsive           & Dashboard usable on both desktop and tablet widths. \\
Offline capability   & Mobile app queues a submitted report locally (SQLite) and sends it once connectivity returns. \\
Performance          & Report submission end-to-end in under 2 minutes; individual API responses target under 2 seconds. \\
Security             & See the dedicated security requirements below --- treated as a first-class, independently-audited concern rather than folded into general performance targets. \\
Compatibility        & iOS 14+, Android 8.0 (API 26)+; modern evergreen browsers for the dashboard. \\
Maintainability / evolution & Automated test suite and CI pipeline (Chapter~\ref{chap:realisation}) so the codebase can evolve without manual regression testing. \\
\bottomrule
\end{tabularx}
\caption{Non-functional requirements}
\label{tab:nfr}
\end{table}

\textbf{Security requirements} deserve their own enumeration, since they were the subject of a dedicated audit phase (Chapter~\ref{chap:realisation}):
\begin{itemize}[leftmargin=*]
  \item All cryptographic signing keys must be rotatable and never retrievable from source control.
  \item Every endpoint that reads or modifies user data must be authenticated and authorized by role.
  \item Sensitive endpoints (login, OTP, password reset) must be rate-limited independently of the platform-wide default.
  \item Sessions must be revocable server-side (logout must have an immediate effect, not just a client-side token discard).
  \item Concurrent requests against the same security-sensitive resource (a refresh token, a one-time ticket) must not both succeed.

\end{itemize}

\subsection{Constraints}

\begin{itemize}[leftmargin=*]
  \item \textbf{Technical}: the stack was fixed early (Flutter, FastAPI, React/TypeScript, PostgreSQL+PostGIS) based on team familiarity and the need for a geospatial-capable relational database; changing any of these mid-internship was not considered.
  \item \textbf{Temporal}: a summer internship is time-boxed to a few months, which shaped the decision to prioritize a working, secure core platform over building the (larger, higher-risk) AI classification microservice.
  \item \textbf{Material}: no dedicated server budget was available at the start, which drove the choice of managed platforms with usable free/low-cost tiers (Render, Vercel, Supabase) over self-hosted infrastructure.
  \item \textbf{Regulatory}: personal data handling needed to be mindful of Tunisian data protection law~\cite{law632004} (Law 63-2004), which informed decisions such as supporting anonymous report submission.
\end{itemize}

\subsection{General Use Case Diagram}

Figure~\ref{fig:usecase} summarizes the main actor--use-case relationships described above.

\begin{figure}[h!]
\centering
\begin{tikzpicture}[
  actor/.style={minimum height=1.6cm},
  usecase/.style={draw=primary, thick, ellipse, minimum width=3.1cm, minimum height=0.9cm, align=center, font=\small},
  boundary/.style={draw=muted, thick, rounded corners=6pt},
]

  % --- actor drawing macro ---
  \newcommand{\stickfigure}[3]{%
    \begin{scope}[shift={(#1,#2)}]
      \draw[thick] (0,0.85) circle (0.18);
      \draw[thick] (0,0.67) -- (0,0);
      \draw[thick] (-0.28,0.4) -- (0.28,0.4);
      \draw[thick] (0,0) -- (-0.22,-0.45);
      \draw[thick] (0,0) -- (0.22,-0.45);
      \node[font=\small\bfseries, below=0.15cm] at (0,-0.45) {#3};
    \end{scope}
  }

  % Actors (left side)
  \stickfigure{-6.3}{2.6}{Citizen}
  \stickfigure{-6.3}{0.6}{Field Agent}
  \stickfigure{-6.3}{-1.4}{Supervisor}

  % Actors (right side)
  \stickfigure{6.3}{2.0}{Analyst}
  \stickfigure{6.3}{-0.6}{Admin}

  % System boundary
  \node[boundary, minimum width=8.6cm, minimum height=6.4cm] (sys) at (0,0.6) {};
  \node[above, font=\bfseries] at (sys.north) {Sahali Platform};

  % Use cases
  \node[usecase] (uc1) at (-2.6,2.6) {Submit Report};
  \node[usecase] (uc2) at (0,2.6)    {Track Report};
  \node[usecase] (uc3) at (2.6,2.6)  {View Statistics};
  \node[usecase] (uc4) at (-2.6,0.6) {Update Status};
  \node[usecase] (uc5) at (0,0.6)    {File Resolution Report};
  \node[usecase] (uc6) at (2.6,0.6)  {Reclassify Report};
  \node[usecase] (uc7) at (-2.6,-1.4){Assign Report};
  \node[usecase] (uc8) at (0,-1.4)   {Manage Team};
  \node[usecase] (uc9) at (2.6,-1.4) {Manage Staff \& Municipalities};

  % Connections
  \draw (-5.5,2.6) -- (uc1.west);
  \draw (-5.5,2.6) -- (uc2.west);
  \draw (-5.5,0.6) -- (uc4.west);
  \draw (-5.5,0.6) -- (uc5.west);
  \draw (-5.5,-1.4) -- (uc7.west);
  \draw (-5.5,-1.4) -- (uc8.west);

  \draw (5.5,2.0) -- (uc3.east);
  \draw (5.5,2.0) -- (uc6.east);
  \draw (5.5,-0.6) -- (uc9.east);
  \draw (5.5,-0.6) -- (uc7.east);

\end{tikzpicture}
\caption{General use case diagram of the Sahali platform}
\label{fig:usecase}
\end{figure}

\chapter{Architecture and Design}
\label{chap:architecture}

\section{General Architecture (Technology-Agnostic)}

At the highest level, Sahali is a \textbf{client--server web and mobile application}: two independent clients (a mobile app and a web dashboard) communicate with a single backend over a REST API, which in turn is the only component with direct access to the database. Figure~\ref{fig:arch-generic} shows this at a technology-agnostic level.

\begin{figure}[h!]
\centering
\begin{tikzpicture}[
  box/.style={draw=primary, thick, rounded corners=4pt, minimum width=3.6cm, minimum height=1cm, align=center, font=\small\bfseries, fill=primary!5},
  arrow/.style={-Latex, thick, color=primary},
]
  \node[box] (mobile) at (-3,3) {Mobile Client\\\footnotesize (Citizens)};
  \node[box] (web) at (3,3) {Web Client\\\footnotesize (Municipal Staff)};
  \node[box, minimum width=5cm] (api) at (0,1) {Backend API};
  \node[box] (db) at (-2.2,-1.2) {Database\\\footnotesize (Geospatial)};
  \node[box] (storage) at (2.2,-1.2) {File Storage\\\footnotesize (Photos, video, audio)};
  \node[box, minimum width=5cm] (ext) at (0,-3.2) {External Services\\\footnotesize (Push, SMS, Email, Geocoding)};

  \draw[arrow] (mobile) -- (api);
  \draw[arrow] (web) -- (api);
  \draw[arrow] (api) -- (db);
  \draw[arrow] (api) -- (storage);
  \draw[arrow] (api) -- (ext);
\end{tikzpicture}
\caption{Technology-agnostic high-level architecture}
\label{fig:arch-generic}
\end{figure}

\textbf{Type of architecture}: layered client--server, not microservices. A single backend process serves both clients; there is no service mesh or independent deployable services beyond the one API (the AI classification component discussed in Chapter~\ref{chap:etude} is architected as a separate microservice but was not implemented during this internship). Background work (notification dispatch, geocoding, municipality matching) runs as in-process asynchronous tasks within the same API process rather than a separate worker fleet.

\textbf{Type of database}: relational, with a geospatial extension. The core entities (users, reports, municipalities, categories) have well-defined relationships that map naturally to a relational schema; the geospatial extension is what specifically distinguishes this from a generic CRUD application, since report locations and municipality boundaries are first-class spatial data, not just latitude/longitude numbers.

\section{Architecture with Technologies}

Figure~\ref{fig:arch-tech} refines the diagram above with the concrete technology choices.

\begin{figure}[h!]
\centering
\begin{tikzpicture}[
  box/.style={draw=primary, thick, rounded corners=4pt, minimum width=4.2cm, minimum height=1.1cm, align=center, font=\footnotesize, fill=primary!5},
  arrow/.style={-Latex, thick, color=primary},
]
  \node[box] (mobile) at (-3.2,3.4) {\textbf{Flutter} (Dart)\\GoRouter, Provider (MVVM)};
  \node[box] (web) at (3.2,3.4) {\textbf{React + TypeScript}\\Vite, Tailwind CSS};
  \node[box, minimum width=6cm] (api) at (0,1.4) {\textbf{FastAPI} (Python 3.12)\\SQLAlchemy 2.0, Alembic, Pydantic v2};
  \node[box] (db) at (-2.6,-1) {\textbf{PostgreSQL 15}\\+ PostGIS};
  \node[box] (redis) at (0.3,-1) {\textbf{Redis}\\OTP, rate limits, SSE pub/sub};
  \node[box] (storage) at (3,-1) {\textbf{Supabase Storage}\\/ MinIO (S3-compatible)};
  \node[box, minimum width=6.5cm] (ext) at (0,-3.4) {\textbf{Firebase FCM} · \textbf{Twilio} · \textbf{SendGrid} · \textbf{OpenStreetMap Nominatim}};

  \draw[arrow] (mobile) -- (api);
  \draw[arrow] (web) -- (api);
  \draw[arrow] (api) -- (db);
  \draw[arrow] (api) -- (redis);
  \draw[arrow] (api) -- (storage);
  \draw[arrow] (api) -- (ext);
\end{tikzpicture}
\caption{Architecture with the selected technology stack}
\label{fig:arch-tech}
\end{figure}

Deployment-wise, the dashboard is a static single-page application hosted on \textbf{Vercel} (deployed on every push to the main branch), the backend runs as a Docker container on \textbf{Render} (running Alembic migrations on startup, then Uvicorn), and \textbf{Supabase} provides the managed production PostgreSQL instance and object storage. Local development uses Docker Compose to run PostgreSQL/PostGIS, Redis, and MinIO.

This stack combines FastAPI~\cite{fastapi} and SQLAlchemy~\cite{sqlalchemy} on the backend, Flutter~\cite{flutter} on mobile, React~\cite{react} on the dashboard, PostGIS~\cite{postgis} for geospatial storage, and Redis~\cite{redis} for caching, rate-limit counters, and pub/sub.

\section{Design}

\subsection{Class Diagram}
\label{sec:class-diagram}

Figure~\ref{fig:class} shows the core domain entities and their relationships, simplified from the full SQLAlchemy schema for readability (junction/audit fields such as timestamps are omitted).

\begin{figure}[h!]
\centering
\resizebox{\textwidth}{!}{%
\begin{tikzpicture}[
  cls/.style={draw=primary, thick, rectangle, align=left, font=\footnotesize, fill=primary!5, inner sep=6pt},
  rel/.style={-Latex, thick},
]
  \node[cls] (user) at (0,6) {
    \begin{tabular}{@{}l@{}}
      \textbf{User} \\ \midrule
      id: UUID \\
      role: citizen\,\textbar\,field\_agent\,\textbar\,analyst\,\textbar\,supervisor\,\textbar\,admin \\
      fullName, phone, email \\
      municipalityId (staff only) \\
      preferredLanguage
    \end{tabular}
  };
  \node[cls] (report) at (7,6) {
    \begin{tabular}{@{}l@{}}
      \textbf{Report} \\ \midrule
      id: UUID, trackingCode \\
      citizenId, categoryId, municipalityId \\
      location: Point (PostGIS) \\
      address, addressAr, city, cityAr \\
      status, priority \\
      aiCategoryId, aiConfidence, isDuplicate
    \end{tabular}
  };
  \node[cls] (municipality) at (0,2.8) {
    \begin{tabular}{@{}l@{}}
      \textbf{Municipality} \\ \midrule
      id, name \\
      location: Point (centroid) \\
      subscriptionTier
    \end{tabular}
  };
  \node[cls] (category) at (7,2.8) {
    \begin{tabular}{@{}l@{}}
      \textbf{Category} \\ \midrule
      id, parentId (self-ref.) \\
      labelAr, labelFr, labelEn \\
      slaHours, isActive
    \end{tabular}
  };
  \node[cls] (assignment) at (0,-0.3) {
    \begin{tabular}{@{}l@{}}
      \textbf{Assignment} \\ \midrule
      id, reportId, agentId \\
      assignedBy, isActive
    \end{tabular}
  };
  \node[cls] (resolution) at (7,-0.3) {
    \begin{tabular}{@{}l@{}}
      \textbf{ResolutionReport} \\ \midrule
      id, reportId, resolvedBy \\
      comment, materials \\
      photoUrls[], videoUrl, voiceNoteUrl
    \end{tabular}
  };
  \node[cls] (history) at (3.5,-3.2) {
    \begin{tabular}{@{}l@{}}
      \textbf{ReportStatusHistory} \\ \midrule
      id, reportId \\
      fromStatus, toStatus, changedBy, note
    \end{tabular}
  };

  \draw[rel] (user) -- node[above, font=\tiny]{1 -- 0..*} (report);
  \draw[rel] (report) -- node[right, font=\tiny]{0..*} (municipality);
  \draw[rel] (report) -- node[left, font=\tiny]{1} (category);
  \draw[rel] (report) -- node[left, font=\tiny]{1 -- 0..*} (assignment);
  \draw[rel] (assignment) -- node[below, font=\tiny]{0..* -- 1} (user);
  \draw[rel] (report) -- node[right, font=\tiny]{1 -- 0..1} (resolution);
  \draw[rel] (report) -- node[below, font=\tiny]{1 -- 0..*} (history);
  \draw[rel] (municipality) -- node[left, font=\tiny]{1 -- 0..*} (user);
\end{tikzpicture}%
}
\caption{Simplified class diagram of the core domain model}
\label{fig:class}
\end{figure}

\subsection{Sequence Diagram --- Report Submission}

Figure~\ref{fig:sequence} traces the full submission flow for a new report, from the mobile client to the background tasks it triggers.

\begin{figure}[h!]
\centering
\begin{tikzpicture}[
  lifeline/.style={dashed, thick, color=muted},
  actor/.style={font=\footnotesize\bfseries, align=center},
  msg/.style={-Latex, thick, color=primary, font=\scriptsize},
]
  \def\ya{9} \def\yb{-6.5}
  \foreach \x/\name in {0/Mobile App, 3/Backend API, 6/PostgreSQL, 9/Nominatim, 12/Notification Services} {
    \node[actor] at (\x,\ya+0.4) {\name};
    \draw[lifeline] (\x,\ya) -- (\x,\yb);
  }

  \draw[msg] (0,8.2) -- (3,8.2) node[midway, above]{POST /reports};
  \draw[msg] (3,7.4) -- (6,7.4) node[midway, above]{INSERT report (status=submitted)};
  \draw[msg] (6,6.6) -- (3,6.6) node[midway, above]{OK};
  \draw[msg] (3,5.8) -- (0,5.8) node[midway, above]{201 Created + tracking code};

  \node[font=\scriptsize\itshape, color=muted] at (6,4.9) {\textit{background tasks, non-blocking}};

  \draw[msg] (3,4.2) -- (6,4.2) node[midway, above]{nearest-municipality query (PostGIS KNN)};
  \draw[msg] (6,3.4) -- (3,3.4) node[midway, above]{municipality\_id};
  \draw[msg] (3,2.6) -- (9,2.6) node[midway, above]{reverse geocode (lat, lng), FR + AR};
  \draw[msg] (9,1.8) -- (3,1.8) node[midway, above]{address, city (FR + AR)};
  \draw[msg] (3,1.0) -- (6,1.0) node[midway, above]{UPDATE report (municipality, address)};
  \draw[msg] (3,0.2) -- (12,0.2) node[midway, above]{notify citizen (push/SMS) + notify staff};
\end{tikzpicture}
\caption{Sequence diagram: report submission and asynchronous enrichment}
\label{fig:sequence}
\end{figure}

Two design decisions are visible in this diagram. First, municipality matching and address geocoding both happen \emph{after} the report is already persisted and acknowledged to the client --- the citizen is never made to wait on an external geocoding API call. Second, geocoding is deliberately queried twice (French, then Arabic) since OpenStreetMap frequently only tags a given point in one language.

\subsection{Package Diagram}

Figure~\ref{fig:package} shows the top-level module organization on both the mobile and backend codebases.

\begin{figure}[h!]
\centering
\begin{tikzpicture}[
  pkg/.style={draw=primary, thick, rectangle, minimum width=3.4cm, minimum height=0.9cm, align=center, font=\footnotesize, fill=primary!5},
  group/.style={draw=muted, thick, rounded corners=4pt, inner sep=10pt},
]
  \begin{scope}
    \node[pkg] (m1) at (0,2) {core\\\footnotesize router, theme, l10n, network};
    \node[pkg] (m2) at (0,0.6) {features\\\footnotesize auth, report, agent, notifications};
    \node[pkg] (m3) at (0,-0.8) {shared\\\footnotesize widgets, models};
    \node[above=0.3cm of m1, font=\bfseries] {Mobile (Flutter)};
  \end{scope}

  \begin{scope}[xshift=6.5cm]
    \node[pkg] (b1) at (0,2) {routers\\\footnotesize auth, reports, admin, events\ldots};
    \node[pkg] (b2) at (0,0.6) {services\\\footnotesize geocoding, storage, otp, notification};
    \node[pkg] (b3) at (0,-0.8) {models / schemas / utils};
    \node[above=0.3cm of b1, font=\bfseries] {Backend (FastAPI)};
  \end{scope}

  \draw[-Latex, thick] (m1) -- (m2);
  \draw[-Latex, thick] (m2) -- (m3);
  \draw[-Latex, thick] (b1) -- (b2);
  \draw[-Latex, thick] (b2) -- (b3);
\end{tikzpicture}
\caption{Package diagram: mobile and backend module organization}
\label{fig:package}
\end{figure}

\chapter{Implementation}
\label{chap:realisation}

\section{Development Environment}

\begin{table}[h!]
\centering
\begin{tabularx}{\textwidth}{lX}
\toprule
\textbf{Category} & \textbf{Tools / Technologies} \\
\midrule
Languages          & Python 3.12, Dart (Flutter SDK $\geq$ 3.11), TypeScript \\
Backend framework  & FastAPI, SQLAlchemy 2.0, Alembic (schema migrations), Pydantic v2 \\
Mobile framework   & Flutter, GoRouter (navigation), Provider (state management, MVVM pattern) \\
Frontend framework & React 19, Vite, Tailwind CSS, react-router-dom \\
Database           & PostgreSQL 15 with the PostGIS geospatial extension \\
Cache / messaging  & Redis (OTP storage, rate-limit counters, token/ticket denylists, SSE pub/sub) \\
Package management & \texttt{uv} (Python), \texttt{pub} (Dart), \texttt{npm} (TypeScript) \\
Containerization    & Docker and Docker Compose (local Postgres/Redis/MinIO) \\
Version control     & Git, hosted on GitHub \\
CI/CD               & GitHub Actions \\
Editor / agent       & Claude Code (LLM-assisted development, see Section~\ref{chap:etude}) \\
\bottomrule
\end{tabularx}
\caption{Development environment summary}
\end{table}

\section{Code Description and Key Algorithms}

This section highlights four representative pieces of the implementation, chosen because each embeds a non-obvious design decision rather than boilerplate CRUD code.

\subsection{Nearest-Municipality Matching}

Rather than requiring a citizen or an operator to pick a municipality manually, every submitted report is matched to the geographically nearest one using a PostGIS $k$-nearest-neighbour query on municipality centroids (Listing~\ref{lst:knn}). This runs synchronously against the database with no external network call, so it adds negligible latency to report submission.

\begin{lstlisting}[language=SQL, caption={Nearest-municipality query (PostGIS KNN operator)}, label={lst:knn}]
SELECT id
FROM municipalities
ORDER BY location <-> ST_SetSRID(ST_MakePoint(:lng, :lat), 4326)
LIMIT 1;
\end{lstlisting}

\subsection{Asymmetric JWT Authentication}

Authentication tokens follow the JSON Web Token standard~\cite{rfc7519} and are signed with RS256 (asymmetric RSA signing) rather than a symmetric HMAC scheme, so that only the backend holding the private key can mint valid tokens, while the public key can be distributed freely for verification. Listing~\ref{lst:jwt} shows token creation, including a \texttt{jti} (JWT ID) claim added specifically to support the revocation mechanism described in Section~\ref{sec:difficultes}.

\begin{lstlisting}[language=Python, caption={Access token creation with a revocation-capable jti claim}, label={lst:jwt}]
def create_access_token(subject: str, role: str) -> str:
    expire = datetime.now(timezone.utc) + timedelta(
        minutes=settings.JWT_ACCESS_TOKEN_EXPIRE_MINUTES
    )
    payload = {
        "sub": subject, "role": role, "exp": expire,
        "type": "access", "jti": str(uuid.uuid4()),
    }
    return jwt.encode(payload, _get_private_key(), algorithm="RS256")
\end{lstlisting}

\subsection{Atomic Single-Use Token Claims}

A recurring theme in the security-hardening phase (Section~\ref{sec:difficultes}) was replacing check-then-act patterns with single atomic Redis~\cite{redis} operations. Listing~\ref{lst:atomic} shows the final refresh-token rotation logic: \texttt{SET\ldots NX} atomically claims a token identifier for one-time use in a single round trip, closing a race window that a separate ``check if used, then mark as used'' sequence would leave open between concurrent requests. This class of vulnerability --- security controls that are correct sequentially but not under concurrency --- is one the OWASP Top 10~\cite{owasp} groups under broken access control and insecure design.

\begin{lstlisting}[language=Python, caption={Atomic one-time-use claim via Redis SET...NX}, label={lst:atomic}]
def consume_once(jti: str, expires_at: datetime) -> bool:
    ttl = _ttl_seconds(expires_at)
    if ttl <= 0:
        return False
    return bool(_redis.set(f"{_PREFIX}{jti}", "1", ex=ttl, nx=True))
\end{lstlisting}

\subsection{Bilingual Reverse Geocoding}

Report addresses are reverse-geocoded through OpenStreetMap's Nominatim service~\cite{nominatim}, queried twice per report --- once with a French language preference, once with Arabic --- because a given point is frequently tagged in only one of the two languages in OpenStreetMap's Tunisia dataset. A conservative resolution strategy discards unreliable neighbourhood-level guesses rather than displaying a plausible-looking but wrong street name (a real issue found during manual testing, discussed in Section~\ref{sec:difficultes}).

\section{Difficulties Encountered and Solutions Adopted}
\label{sec:difficultes}

Table~\ref{tab:difficulties} summarizes the most significant problems encountered during the internship and how each was resolved. They are grouped by the project phase in which they arose.

\begin{table}[h!]
\centering
\begin{tabularx}{\textwidth}{p{3.6cm}Xp{4.7cm}}
\toprule
\textbf{Difficulty} & \textbf{Root Cause} & \textbf{Solution} \\
\midrule
Field agents received ``report not found'' for reports they were legitimately assigned to. &
A municipality-scoping access check, added for admin/supervisor roles, was applied uniformly to field agents too --- but a report's \texttt{municipality\_id} was often \texttt{NULL} (not yet backfilled), and \texttt{None != <int>} evaluates to \texttt{True} in Python, so the comparison always failed. &
Gave field agents their own access-check branch based on the \texttt{Assignment} table (task-based access) instead of municipality comparison (membership-based access) --- the two roles have fundamentally different access models. \\
\addlinespace
Reverse-geocoded street names were sometimes plainly wrong (a neighbourhood name shown as a street). &
Confirmed, by testing across every Nominatim zoom level, that this was a genuine OpenStreetMap data-quality gap for the queried point, not a bug in the query. &
Removed the risky suburb/neighbourhood fallback; a report is left without a street name rather than shown an incorrect one. \\
\addlinespace
Android release build failed with an abstract-class implementation error in a third-party package. &
A version-incompatible pairing between the \texttt{record} plugin and its platform-interface package, resolved incorrectly by the dependency resolver at an older pinned version. &
Bumped the \texttt{record} constraint to a version range with a compatible pairing; confirmed via a clean static analysis pass and a full release build. \\
\addlinespace
A signing key was found retrievable from an early Git commit, still reachable from every branch. &
The key had been committed once, early in the project's history, and never removed --- deleting it from the current checkout does not remove it from history. &
Rotated the key pair immediately (treating the old one as permanently compromised), then separately evaluated (and deferred, as a lower-priority hygiene step) rewriting Git history to purge the blob. \\
\addlinespace
The CI dependency-vulnerability scan failed with dozens of known CVEs once it could finally reach the internet (a local machine issue had silently prevented it from running at all until then). &
Several pinned dependencies were multiple minor versions behind, and one framework version capped how far a related package could be upgraded, requiring the upgrade chain to be resolved in the right order. &
Methodically bisected compatible version combinations by installing candidates into a disposable virtual environment and re-running the full test suite after each change, rather than upgrading blindly. \\
\addlinespace
An independent automated code review found that two ``fixed'' security issues (token rotation, one-time ticket consumption) still had a narrow race window. &
The fixes had replaced an unsafe pattern with a \emph{check, then act} pattern --- safer than the original, but still not atomic under concurrent requests. &
Replaced both with single atomic Redis operations (\texttt{SET...NX}, \texttt{GETDEL}); added a test using a synchronization barrier to force two threads to race deliberately, rather than hoping a naive concurrent test would happen to overlap. \\
\addlinespace
A Gradle check added to fail loudly when the release-signing keystore was missing ended up breaking \emph{all} builds, including debug ones. &
The check ran at Gradle's configuration phase, which evaluates for every invocation, not only when a release build is actually requested. &
Moved the check to fire immediately before the actual signing task instead, verified by testing all three cases: no keystore with a debug build (succeeds), no keystore with a release build (fails with a clear message), and a real keystore with a release build (produces a valid signed artifact). \\
\bottomrule
\end{tabularx}
\caption{Significant difficulties encountered and their resolutions}
\label{tab:difficulties}
\end{table}

A pattern across several of these entries is worth naming explicitly: more than one ``fix'' turned out to be incomplete on first attempt, and was only caught by a subsequent, independent verification step (a fresh test, an automated review pass, or a deliberately adversarial concurrency test). This is the practical justification for the audit/fix/verify/review methodology described in Section~\ref{chap:cadre-general} --- verification was not a formality, it repeatedly changed the outcome.

\section{User Interfaces}

Figures~\ref{fig:ui-mobile} and~\ref{fig:ui-dashboard} illustrate the citizen-facing mobile flow and the municipal staff dashboard respectively.

\begin{figure}[h!]
\centering
\fbox{\parbox{0.85\textwidth}{\centering\vspace{3cm}
\textit{[Insert screenshots here: onboarding / language selection, report submission wizard (category → photo → location → description → review), report tracking screen, field-agent mission list]}
\vspace{3cm}}}
\caption{Mobile application --- key citizen and field-agent screens}
\label{fig:ui-mobile}
\end{figure}

\begin{figure}[h!]
\centering
\fbox{\parbox{0.85\textwidth}{\centering\vspace{3cm}
\textit{[Insert screenshots here: dashboard overview with KPIs, live report map, report detail panel, team/agent management view]}
\vspace{3cm}}}
\caption{Web dashboard --- key municipal staff screens}
\label{fig:ui-dashboard}
\end{figure}

\section{CI/CD Pipeline}
\label{sec:cicd}

Once the functional platform and the security-hardening pass were in place, a continuous integration and deployment pipeline was built using GitHub Actions~\cite{gha}, structured as four independent jobs on every pull request:

\begin{enumerate}[leftmargin=*]
  \item \textbf{Backend lint and dependency audit} --- static analysis (ruff) and a known-vulnerability scan (pip-audit) against the pinned dependency set.
  \item \textbf{Backend test suite} --- runs the automated test suite (Chapter~\ref{chap:tests}) against real PostgreSQL and Redis service containers, mirroring production dependencies rather than mocking them.
  \item \textbf{Backend Docker build check} --- builds the production container image to catch a broken \texttt{Dockerfile} or dependency-installation failure before the hosting platform does.
  \item \textbf{Dashboard and mobile checks} --- type-checking and build for the dashboard, static analysis and the mobile test suite for the Flutter app.
\end{enumerate}

A separate, tag-triggered workflow automates mobile releases: pushing a version tag builds a signed Android APK and app bundle and attaches them to a GitHub Release, using the signing keystore and its passwords stored as encrypted repository secrets --- never committed to source control. Because this workflow holds write access to the repository and to those secrets, every third-party action it uses is pinned to an immutable commit hash rather than a mutable version tag, closing off a supply-chain risk where a compromised action release could otherwise run arbitrary code with those permissions.

\chapter{Testing and Validation}
\label{chap:tests}

\section{Testing Methodology}

Two complementary layers of testing were used throughout the project.

\textbf{Manual, scripted verification} was used continuously during the security-hardening phase: for every fix, a direct sequence of API calls (via \texttt{curl}) was run against a live local instance of the backend to confirm the exact before/after behavior change --- for example, sending the same login request eleven times in under a minute to confirm the twelfth was rejected with \texttt{429 Too Many Requests}, or reusing a refresh token twice to confirm the second use was rejected. This caught issues an automated test alone would not have surfaced as quickly, since it exercised the real HTTP layer, the real Redis instance, and the real rate-limiter middleware together.

\textbf{Automated regression testing} was added once the platform's behavior stabilized, as part of the CI/CD pipeline (Section~\ref{sec:cicd}). The backend test suite runs against real PostgreSQL and Redis service containers in CI (not mocks), on the reasoning that mocking the exact services being tested for atomicity and timing behavior would test the mock, not the system. Concurrency-sensitive behavior (Section~\ref{sec:difficultes}) is tested with a \texttt{threading.Barrier} that forces two consumer threads to race at the same instant, rather than a plain thread pool that does not guarantee the two calls actually overlap.

\section{Test Cases and Results}

Table~\ref{tab:testcases} lists the automated backend test suite, all of which pass in the CI pipeline.

\begin{table}[h!]
\centering
\begin{tabularx}{\textwidth}{Xl}
\toprule
\textbf{Test case} & \textbf{Result} \\
\midrule
OTP verification locks out after five wrong attempts, rejecting even the correct code afterward & Pass \\
A fresh OTP request resets the attempt lockout & Pass \\
A refresh token can be used exactly once; a second use is rejected & Pass \\
Logout immediately revokes both the access and refresh token & Pass \\
The login endpoint's rate limit rejects the 11th request within a minute & Pass \\
SSE event-stream ticket issuance requires a staff role (a citizen is rejected) & Pass \\
An SSE ticket can be consumed exactly once & Pass \\
Two threads racing to consume the same SSE ticket --- exactly one succeeds & Pass \\
An invalid/unknown SSE ticket is rejected & Pass \\
\bottomrule
\end{tabularx}
\caption{Automated backend test cases and results}
\label{tab:testcases}
\end{table}

Beyond this automated suite, the manual verification pass covered: end-to-end registration and login (including real email validation), the full report submission and geocoding flow, Android release build signing under both success and failure conditions (Section~\ref{sec:difficultes}), and a from-scratch dependency install (a disposable virtual environment populated only from the pinned requirements file) to confirm the pinned dependency set is self-consistent and not merely working by accident of an already-populated development environment.

\section{Performance Analysis}

Formal load testing under realistic concurrent traffic was \textbf{not performed} during this internship and is recorded here as a limitation rather than omitted silently. What was measured directly:

\begin{itemize}[leftmargin=*]
  \item Individual API response times observed during manual testing were consistently under the 2-second target for all tested endpoints on local infrastructure.
  \item The rate limiter was confirmed to enforce its configured limit correctly \emph{across} requests sharing Redis-backed storage (as opposed to per-process in-memory counters, which would under-enforce the limit across multiple backend worker processes) --- itself a performance-adjacent correctness issue caught during the audit, not a raw throughput measurement.
\end{itemize}

Load testing against the deployed environment (concurrent report submissions, dashboard SSE connections under load) is listed as future work in Chapter~\ref{chap:conclusion}.

\section{Comparison with the Original Non-Functional Objectives}

Table~\ref{tab:nfr-status} revisits the non-functional requirements from Section~\ref{sec:etat-art} and states, honestly, which were verified, which were partially addressed, and which remain open.

\begin{table}[h!]
\centering
\begin{tabularx}{\textwidth}{lXl}
\toprule
\textbf{Requirement} & \textbf{Status} & \textbf{Note} \\
\midrule
Multilingual         & Achieved & Mobile: AR/FR/EN with RTL. Dashboard: FR/AR. \\
Offline capability   & Achieved & Local SQLite queue on the mobile app, flushed on reconnect. \\
Security             & Achieved & Three-phase audit, independently re-verified by automated code review. \\
Performance ($<$2s)  & Partially verified & Confirmed manually; not load-tested (see above). \\
Compatibility        & Partially verified & Verified for the primary target versions; not tested across the full device/browser matrix. \\
Maintainability / evolution & Achieved & Automated test suite + CI pipeline in place. \\
Responsive dashboard & Not formally verified & No dedicated cross-device layout testing was performed. \\
\bottomrule
\end{tabularx}
\caption{Non-functional requirement verification status}
\label{tab:nfr-status}
\end{table}

\chapter{Conclusion and Perspectives}
\label{chap:conclusion}

\section{Project Assessment}

This internship set out to build a working, secure civic reporting platform for Tunisian municipalities, and delivered exactly that: a trilingual mobile application, a municipal staff dashboard, and a backend API, all communicating over a REST interface backed by a geospatial relational database. Beyond the functional core, two additional outcomes turned out to be at least as significant as the original scope: a full, severity-ranked security audit that found and fixed a genuinely exploitable set of issues (a retrievable signing key, unauthenticated sensitive endpoints, non-revocable sessions), and a continuous integration/deployment pipeline that itself surfaced real, previously-undetected defects --- a mobile smoke test that had silently never passed, and several dependencies with known vulnerabilities that had drifted unnoticed.

If there is one overall lesson from the internship, it is that \textbf{verification changes outcomes, not just confidence}: more than one fix in Chapter~\ref{chap:realisation} was incomplete on its first attempt and was only caught by a subsequent, independent check. Treating audit, fix, and verification as separate, repeated steps --- rather than folding verification into ``I fixed it, moving on'' --- was the single methodological choice that mattered most.

\section{Personal and Technical Contributions}

On the technical side, this internship required working across the full stack of a real product: a Flutter mobile application with offline support and trilingual RTL-aware UI, a React/TypeScript dashboard, a FastAPI backend with a PostGIS-backed geospatial data model, and a GitHub Actions CI/CD pipeline with real service-container-based integration tests. It also required a security-specific skill set distinct from feature development: reasoning about attacker-visible attack surface (what can an unauthenticated request reach?), about race conditions in concurrent request handling, and about the operational side of key rotation and credential hygiene in a deployed system, not just a textbook description of them.

On the personal side, working with an LLM-assisted development workflow (Section~\ref{chap:etude}) required developing a specific discipline: treating the agent's output as a proposal to verify, not a finished answer to trust, and building that verification into the workflow itself (a curl-based regression check, an automated test, an independent review pass) rather than relying on visual code inspection alone.

\section{Possible Improvements and Perspectives}

The following items were identified during the internship as genuine next steps, not filler:

\begin{itemize}[leftmargin=*]
  \item \textbf{Municipality-scoped staff visibility} --- reports now auto-populate their nearest municipality, but non-admin staff roles still see every report platform-wide rather than only their own municipality's. This is the most immediately impactful next feature.
  \item \textbf{AI classification microservice} --- the data model already reserves fields for AI-suggested category, confidence score, and duplicate detection; the model itself was out of scope for this internship and remains to be built.
  \item \textbf{Load testing} --- formal concurrent-load testing against the deployed environment, to validate the sub-2-second response target under realistic traffic rather than only during manual single-user testing.
  \item \textbf{Git history remediation} --- purging the historical commit containing the now-rotated signing key from the repository's history entirely, deferred as a lower-urgency hygiene task once the key itself was rotated and the repository was made private.
  \item \textbf{Upload content validation} --- the current upload path validates a file's declared content type; validating the actual file signature (``magic bytes'') server-side would close a narrow gap where a client could mislabel an upload's type.
  \item \textbf{Email verification persistence} --- the email-verification endpoint currently confirms a one-time code but does not yet persist a verified/unverified state on the account, since nothing in the current product gates behavior on it; this is a product decision to make before it becomes a schema change.
\end{itemize}

Taken together, these are consistent with a platform that reached a genuinely working, secure, and continuously-tested state within the internship's timeframe, with a clear and specific --- rather than vague --- list of what comes next.


% ============================================================
% ANNEXES
% ============================================================
\appendix
\chapter{Report Category Taxonomy and SLA Targets}

\begin{table}[h!]
\centering
\begin{tabularx}{\textwidth}{Xl}
\toprule
\textbf{Category} & \textbf{Target Resolution Time} \\
\midrule
Safety and Security                          & 24 hours \\
Public Lighting --- exposed wiring           & 24 hours \\
Water and Sanitation                         & 48 hours \\
Public Lighting --- faulty streetlights      & 72 hours \\
Infrastructure (potholes, sidewalks, signage) & 7 days \\
Cleanliness and Waste Management             & 7 days \\
Transportation (bus stops, traffic signals)  & 7 days \\
Environment (pollution, illegal tree cutting) & 14 days \\
\bottomrule
\end{tabularx}
\caption{Default service-level targets by report category}
\end{table}

\chapter{CI Pipeline Excerpt}

Listing~\ref{lst:ci} shows the backend test job from the GitHub Actions workflow, including the real PostgreSQL/PostGIS and Redis service containers used in CI.

\begin{lstlisting}[language=bash, caption={Backend test job (excerpt), \texttt{.github/workflows/ci.yml}}, label={lst:ci}]
backend-tests:
  runs-on: ubuntu-latest
  services:
    postgres:
      image: postgis/postgis:15-3.4
      env:
        POSTGRES_USER: citizen_alert
        POSTGRES_PASSWORD: password
        POSTGRES_DB: citizen_alert_test
      ports: ["5432:5432"]
    redis:
      image: redis:7-alpine
      ports: ["6379:6379"]
  steps:
    - uses: actions/checkout@v4
      with:
        persist-credentials: false
    - uses: actions/setup-python@v5
      with:
        python-version: "3.12"
    - run: pip install -r requirements.txt pytest==8.2.2
    - run: alembic upgrade head > /dev/null
    - run: python -m pytest -v
\end{lstlisting}

\chapter{Security Audit Summary}

Table~\ref{tab:security-summary} summarizes the audit's findings by severity, as documented in the project's internal \texttt{SECURITY.md}.

\begin{table}[h!]
\centering
\begin{tabularx}{\textwidth}{lXl}
\toprule
\textbf{Severity} & \textbf{Example finding} & \textbf{Status} \\
\midrule
Critical & RS256 signing key retrievable from Git history & Rotated \\
Critical & Photo upload endpoint unauthenticated, no size/type limits & Fixed \\
Critical & No server-side session revocation & Fixed (Redis denylist) \\
High     & CORS accepted any \texttt{*.vercel.app} origin & Partially fixed \\
High     & No per-route rate limiting on sensitive endpoints & Fixed \\
High     & OTP comparison not constant-time, no attempt cap & Fixed \\
Medium   & No CI/automated dependency scanning & Fixed \\
Low      & Non-functional cleartext-traffic allowlist entry & Fixed \\
\bottomrule
\end{tabularx}
\caption{Security audit findings summary}
\label{tab:security-summary}
\end{table}


% ============================================================
% BIBLIOGRAPHIE
% ============================================================
\begin{thebibliography}{99}

\bibitem{fastapi}
S.\ Ramírez, \textit{FastAPI Documentation}. [Online]. Available: \url{https://fastapi.tiangolo.com/}

\bibitem{flutter}
Google, \textit{Flutter Documentation}. [Online]. Available: \url{https://docs.flutter.dev/}

\bibitem{react}
Meta Open Source, \textit{React Documentation}. [Online]. Available: \url{https://react.dev/}

\bibitem{postgis}
PostGIS Project Steering Committee, \textit{PostGIS 3.4 Manual}. [Online]. Available: \url{https://postgis.net/documentation/}

\bibitem{sqlalchemy}
M.\ Bayer, \textit{SQLAlchemy Documentation}, version 2.0. [Online]. Available: \url{https://docs.sqlalchemy.org/}

\bibitem{redis}
Redis Ltd., \textit{Redis Documentation}. [Online]. Available: \url{https://redis.io/docs/}

\bibitem{rfc7519}
M.\ Jones, J.\ Bradley, and N.\ Sakimura, ``JSON Web Token (JWT),'' \textit{IETF RFC 7519}, May 2015. [Online]. Available: \url{https://www.rfc-editor.org/rfc/rfc7519}

\bibitem{owasp}
OWASP Foundation, \textit{OWASP Top 10:2021 --- The Ten Most Critical Web Application Security Risks}. [Online]. Available: \url{https://owasp.org/Top10/}

\bibitem{gha}
GitHub Inc., \textit{GitHub Actions Documentation}. [Online]. Available: \url{https://docs.github.com/en/actions}

\bibitem{fixmystreet}
mySociety, \textit{FixMyStreet --- Open Source Street Reporting Platform}. [Online]. Available: \url{https://www.fixmystreet.com/}

\bibitem{nyc311}
City of New York, \textit{NYC 311 Service Requests}. [Online]. Available: \url{https://portal.311.nyc.gov/}

\bibitem{nominatim}
OpenStreetMap Foundation, \textit{Nominatim --- Search and Reverse Geocoding for OpenStreetMap Data}. [Online]. Available: \url{https://nominatim.org/}

\bibitem{law632004}
Republic of Tunisia, \textit{Loi organique n\textdegree~63-2004 du 27 juillet 2004, portant sur la protection des données à caractère personnel}, Journal Officiel de la République Tunisienne, 2004.

\end{thebibliography}


\end{document}
